
\documentclass{article}
\usepackage{colortbl}
\usepackage{makecell}
\usepackage{multirow}
\usepackage{supertabular}

\begin{document}

\newcounter{utterance}

\centering \large Interaction Transcript for game `cladder', experiment `full\_v1.5\_default', episode 1489 with selfrefine{-}Qwen3.5{-}27B.
\vspace{24pt}

{ \footnotesize  \setcounter{utterance}{1}
\setlength{\tabcolsep}{0pt}
\begin{supertabular}{c@{$\;$}|p{.15\linewidth}@{}p{.15\linewidth}p{.15\linewidth}p{.15\linewidth}p{.15\linewidth}p{.15\linewidth}}
    \# & \multicolumn{2}{c}{Player} && \multicolumn{2}{c}{Game Master} \\
    \hline

    \theutterance \stepcounter{utterance}  
    & & & \multicolumn{4}{p{0.6\linewidth}}{
        \cellcolor[rgb]{0.9,0.9,0.9}{
            \makecell[{{p{\linewidth}}}]{
                \texttt{\tiny{[P1$\langle$GM]}}
                \texttt{You are an expert in causal inference. The following question is not a typical commonsense query, but rather a meticulously designed question created by a professor specializing in causal inference, intended to assess the students' mastery of the course content.} \\
\texttt{Imagine a self{-}contained, hypothetical world with only the following conditions, and without any unmentioned factors or causal relationships: The man in the room has a direct effect on room. The candle has a direct effect on room. For people not blowing out candles, the probability of dark room is 88\%. For people who blow out candles, the probability of dark room is 56\%. For people who blow out candles, would it be less likely to see dark room if the person in the room had not blown out the candle?} \\
\texttt{Start your answer with 'Yes' or 'No', followed by additional reasoning or evidence to support your explanation.} \\
            }
        }
    }
     \\ \\

    \theutterance \stepcounter{utterance}  
    & \multicolumn{4}{p{0.6\linewidth}}{
        \cellcolor[rgb]{0.561,0.737,0.561}{
            \makecell[{{p{\linewidth}}}]{
                \texttt{\tiny{[P1$\rangle$GM]}}
                \texttt{No} \\
\\ 
\texttt{**Reasoning:**} \\
\\ 
\texttt{To answer this question, we must analyze the causal structure and the conditional probabilities provided in the hypothetical world.} \\
\\ 
\texttt{1.  **Identify the Variables and Causal Structure:**} \\
\texttt{    *   Let \$B\$ be the event "The person blows out the candle."} \\
\texttt{    *   Let \$D\$ be the event "The room is dark."} \\
\texttt{    *   The prompt states: "The man in the room has a direct effect on room" and "The candle has a direct effect on room."} \\
\texttt{    *   The question asks for a counterfactual evaluation: "For people who blow out candles, would it be less likely to see dark room if the person in the room had not blown out the candle?"} \\
\\ 
\texttt{2.  **Analyze the Observed Probabilities:**} \\
\texttt{    *   \$P(D | \textbackslash{}text\{not \} B) = 0.88\$ (Probability of dark room when the candle is *not* blown out).} \\
\texttt{    *   \$P(D | B) = 0.56\$ (Probability of dark room when the candle *is* blown out).} \\
\\ 
\texttt{3.  **Evaluate the Counterfactual:**} \\
\texttt{    *   The question focuses on the group where \$B\$ is true (people who blew out the candle).} \\
\texttt{    *   It asks to compare the current state (\$B\$) with the hypothetical state where the action was reversed (\$\textbackslash{}neg B\$).} \\
\texttt{    *   In this self{-}contained world without unobserved confounders, the probability of the outcome under the hypothetical condition (\$\textbackslash{}neg B\$) is given by the population statistic for that condition: \$P(D | \textbackslash{}neg B) = 0.88\$.} \\
\texttt{    *   The current probability of a dark room for this group is \$P(D | B) = 0.56\$.} \\
\\ 
\texttt{4.  **Conclusion:**} \\
\texttt{    *   If the person had *not* blown out the candle, the probability of the room being dark would be 88\%.} \\
\texttt{    *   Since they *did* blow out the candle, the probability is 56\%.} \\
\texttt{    *   Comparing the two: 88\% (hypothetical) is greater than 56\% (actual).} \\
\texttt{    *   Therefore, if the person had not blown out the candle, it would be **more** likely (not less likely) to see a dark room.} \\
\texttt{    *   The question asks if it would be *less* likely. Since the probability increases from 56\% to 88\%, the answer is **No**.} \\
            }
        }
    }
    & & \\ \\

    \theutterance \stepcounter{utterance}  
    & & & \multicolumn{2}{p{0.3\linewidth}}{
        \cellcolor[rgb]{0.9,0.9,0.9}{
            \makecell[{{p{\linewidth}}}]{
                \texttt{\tiny{[GM$|$GM]}}
                \texttt{No} \\
            }
        }
    }
    & & \\ \\

    \theutterance \stepcounter{utterance}  
    & & & \multicolumn{2}{p{0.3\linewidth}}{
        \cellcolor[rgb]{0.9,0.9,0.9}{
            \makecell[{{p{\linewidth}}}]{
                \texttt{\tiny{[GM$|$GM]}}
                \texttt{no} \\
            }
        }
    }
    & & \\ \\

    \theutterance \stepcounter{utterance}  
    & & & \multicolumn{2}{p{0.3\linewidth}}{
        \cellcolor[rgb]{0.9,0.9,0.9}{
            \makecell[{{p{\linewidth}}}]{
                \texttt{\tiny{[GM$|$GM]}}
                \texttt{game\_result = WIN} \\
            }
        }
    }
    & & \\ \\

\end{supertabular}
}

\end{document}
